\section{Related Work}
\label{sec:related}

\subsection{Addressing Temporal Covariate Shift}
The TCS problem was first identified in BNTT~\cite{bntt:2020} and later addressed by methods such as TEBN~\cite{tebn:2022} and TAB~\cite{tab:2024}, which adapt normalization layers on a per-timestep basis. While these methods alleviate TCS in activations, they increase model complexity with additional parameters and, more critically, overlook the persistent shift in membrane potential, a primary source of TCS during inference (see Section~\ref{sec:root}). 

IMP+LTS~\cite{shen:2024} proposed training the initial membrane potential for every neuron to enhance representational power. However, it does not directly address TCS and incurs a large parameter overhead scaling with the number of neurons $O(N)$. By contrast, our approach initializes membrane potentials at the layer level $O(L)$ ($L \ll N$), directly targeting TCS while remaining lightweight.

More recently, MPS~\cite{ding:2025} reinterpreted SNNs from an ensemble-learning perspective, arguing that excessive differences in membrane potential distributions across timesteps destabilize subnetworks. Their smoothing and guidance strategies improve temporal consistency and yield accuracy gains. However, these methods modify the LIF update rule itself and require additional smoothing coefficients and guidance losses, meaning both training and inference diverge from the standard LIF neurons. In contrast, our approach preserves the original inference pipeline: MP-Init adds negligible overhead during training while inference remains compatible with neuromorphic hardware~\cite{davies:2018}.

%Crucially, none of these approaches directly address the TCS in the membrane potential, which is the root cause of TCS in activations, as will be explained in \cref{sec:root}. In contrast, our methods directly address this issue, fundamentally reducing TCS without adding significant overhead, thereby preserving model simplicity and biological plausibility.

\subsection{Learning in Spiking Neural Networks}
Surrogate Gradient (SG) descent enables end-to-end training of SNNs by approximating the non-differentiable firing process~\cite{lee:2016, wu:2018, wu:2019}. Some studies primarily focus on the impact of hyperparameters and the shape of surrogate derivatives~\cite{zenke:2021, li:2021}, while recent advancements introduce differentiable surrogate gradient functions~\cite{li:2021} and Temporal Efficient Training (TET) losses~\cite{deng:2022}. Other approaches regulate membrane potential by adding loss terms~\cite{guo:2022, guo:2023} or applying batch normalization~\cite{guo:2023_2}. Meanwhile, methods that modify the dynamics of LIF neurons~\cite{fang:2023, yao:2022, huang:2024} have been explored, though they could be incompatible with existing neuromorphic hardware~\cite{hunsberger:2015, review_guo:2023}.

Recent studies have also explored training internal parameters, such as the threshold voltage and leakage term, via gradient descent~\cite{fang:2021, wang:2022, rathi:2023}. For instance, PLIF~\cite{fang:2021} optimizes the leakage term, LTMD~\cite{wang:2022} adjusts the threshold voltage, and DIET-SNN~\cite{rathi:2023} learns both. However, these works provide limited analysis of gradient flow related to these learnable parameters, leaving questions about training stability unanswered. In contrast, we find that when the threshold voltage becomes trainable, surrogate gradient methods can exhibit instability, leading to misaligned, exploding, or vanishing gradients, as discussed in Section~\ref{sec:type1_type2}. To overcome these issues, we propose a novel SG formulation that remains stable when learning the threshold voltage, ensuring robust convergence.  
